\section{Shifted harmonic sums, beta derivatives, and exact Laurent data}
\label{harm:sec:main}

The all-order Stieltjes differentiation formulas in the manuscript have
elementary symmetric harmonic coefficients.  Those coefficients can also
be put inside a Dirichlet series.  This gives a second exact calculus:
positive integer values are beta derivatives, whereas all principal parts
and finite parts at nonpositive integers are finite polynomials.  The
dependence on an arbitrary shift is retained throughout.

There are important antecedents.  Young's parametric hyperharmonic Euler
sums include the beta-derivative evaluation below
\cite{Young2026}; his work on height-one multiple zeta functions supplies
the unshifted generating-function framework \cite{Young2024}, and his
earlier paper treats their unshifted Laurent expansions \cite{Young2023}.  Harmonic
Stieltjes constants at depth one have a substantial separate literature
\cite{HarmonicStieltjes2023}.  Accordingly, neither the beta method nor the
unshifted height-one theory is asserted to be new.  The main continuation
developed here is the all-depth coefficient rule for outer-shifted
Laurent data at nonpositive integers, together with its polynomial shift
calculus and the residue term required when finite parts are
differentiated.  These formulas were not located in the inspected
repository sources; no claim of priority over all literature is made.

\subsection{Definitions and a single two-variable family}

Put $A=a+1$ and initially assume $\Re a>-1$.  Write
\begin{equation}\label{harm:eq:elementary}
 e_r(N)=[u^r]\prod_{j=1}^{N}\left(1+\frac uj\right),
 \qquad e_0(N)=1,
 \qquad e_r(N)=0\quad(r>N).
\end{equation}
Thus $e_r(N)=H_N(\{1\}^r)$ is the strict multiple harmonic sum, and
\begin{align*}
 e_1(N)&=H_N,\\
 e_2(N)&=\frac{H_N^2-H_N^{(2)}}2,\\
 e_3(N)&=\frac{H_N^3-3H_NH_N^{(2)}+2H_N^{(3)}}6.
\end{align*}
The main family is
\begin{equation}\label{harm:eq:E}
 E_r(s;a)=\sum_{n=1}^{\infty}\frac{e_r(n-1)}{(n+a)^s},
 \qquad r\geq0,\quad \Re s>1.
\end{equation}
The power uses the principal logarithm; $n+a$ lies in the right half-plane.
For $a=0$, this is $\zeta(s,\{1\}^r)$, and for $r=0$ it is
$\zeta(s,A)$.  The shift affects only the outer denominator.  In
particular, for general $a$ this is not the harmonic Hurwitz function
in which every inner denominator is shifted as well.

The bound $e_r(n-1)\leq H_{n-1}^r/r!$ proves normal convergence of
\eqref{harm:eq:E} on compact subsets of the displayed domain.  To keep
all depths together, define
\begin{equation}\label{harm:eq:F}
 F(s,u;a)=\sum_{n=1}^{\infty}
 \frac{\Gamma(n+u)}{\Gamma(n)\Gamma(1+u)}\frac1{(n+a)^s}.
\end{equation}
For $u$ near zero and $\Re s>1+|u|$, the defining series is normally
convergent.  The quotient in the numerator is the polynomial
$\prod_{j=1}^{n-1}(1+u/j)$, so
\begin{equation}\label{harm:eq:coeff}
 F(s,u;a)=\sum_{r\geq0}E_r(s;a)u^r.
\end{equation}
Equivalently, $F$ is the Barnes--Hurwitz zeta function with complex
order $1+u$, in the convention used in \cite{Young2024}.

\begin{proposition}[Mellin representations]\label{harm:prop:mellin}
In a common domain of absolute convergence,
\begin{align}
 \Gamma(s)F(s,u;a)
 &=\int_0^{\infty}t^{s-1}e^{-At}
                  (1-e^{-t})^{-1-u}\,dt,\label{harm:eq:mellinF}\\
 E_r(s;a)
 &=\frac1{r!\Gamma(s)}\int_0^1
       \frac{x^a(-\log x)^{s-1}[-\log(1-x)]^r}{1-x}\,dx.
       \label{harm:eq:mellinE}
\end{align}
For fixed $r$, the second integral converges for $\Re s>1$ and
$\Re a>-1$.
\end{proposition}
\begin{proof}
The binomial series gives
\[
 \sum_{n\geq1}\frac{\Gamma(n+u)}{\Gamma(n)\Gamma(1+u)}x^n
       =\frac{x}{(1-x)^{1+u}}.
\]
Multiply by $x^{a-1}(-\log x)^{s-1}$ and integrate.  Absolute
convergence permits termwise integration, and the integral of the
$n$th term is $\Gamma(s)(n+a)^{-s}$.  The substitution $x=e^{-t}$
gives \eqref{harm:eq:mellinF}; coefficient extraction in $u$ gives
\eqref{harm:eq:mellinE}.  At $x=1$, its absolute integrand is bounded
by a constant times
$(1-x)^{\Re s-2}(1+|\log(1-x)|^r)$, which is integrable for
$\Re s>1$.  At zero, $\Re a>-1$ suffices.
\end{proof}

This also gives a direct polylogarithmic interpretation, since
\[
 \operatorname{Li}_{\{1\}^{r+1}}(x)
 =\frac{[-\log(1-x)]^{r+1}}{(r+1)!},\qquad
 \frac{d}{dx}\operatorname{Li}_{\{1\}^{r+1}}(x)
 =\frac{[-\log(1-x)]^r}{r!(1-x)}.
\]
Thus \eqref{harm:eq:mellinE} is the Mellin transform of a derivative
of a multiple polylogarithm of arbitrary depth.

\subsection{Positive integer values and their antiderivatives}

Define generalized harmonic functions by
\begin{equation}\label{harm:eq:Hcomplex}
 H_a^{(1)}=\psi(A)+\gamma,\qquad
 H_a^{(k)}=\zeta(k)-\zeta(k,A)\quad(k\geq2).
\end{equation}
They agree with the usual finite harmonic numbers at nonnegative
integer $a$.  Introduce
\begin{equation}\label{harm:eq:R}
 R_a(u)=\frac{\Gamma(1-u)\Gamma(A)}{\Gamma(A-u)}
       =\exp\left(\sum_{k\geq1}\frac{H_a^{(k)}}k u^k\right)
       =\sum_{d\geq0}h_d(a)u^d.
\end{equation}
Here $h_d(a)$ is a complete symmetric harmonic polynomial:
\begin{equation}\label{harm:eq:hrec}
 h_0(a)=1,\qquad
 d\,h_d(a)=\sum_{k=1}^d H_a^{(k)}h_{d-k}(a).
\end{equation}
For example,
\[
 h_1=H_a^{(1)},\quad
 h_2=\frac{(H_a^{(1)})^2+H_a^{(2)}}2,\quad
 h_3=\frac{(H_a^{(1)})^3+3H_a^{(1)}H_a^{(2)}+2H_a^{(3)}}6.
\]
The plus signs distinguish these polynomials from the elementary
coefficients $e_r(N)$ in the summand.

\begin{theorem}[Subtracted beta identity and all positive integer values]
\label{harm:thm:beta}
For $r\geq0$ and $\Re a>-1$,
\begin{equation}\label{harm:eq:primitive}
 \boxed{\quad
 h_{r+1}(a)=\sum_{n\geq1}e_r(n-1)
           \left(\frac1n-\frac1{n+a}\right).
 \quad}
\end{equation}
For every integer $p\geq2$,
\begin{equation}\label{harm:eq:positive}
 \boxed{\quad
 E_r(p;a)=\frac{(-1)^p}{(p-1)!}
               \frac{d^{p-1}}{da^{p-1}}h_{r+1}(a).
 \quad}
\end{equation}
Equivalently, as an analytic germ at $(x,u)=(0,0)$,
\begin{equation}\label{harm:eq:betaG}
 \sum_{p\geq2}\sum_{r\geq0}E_r(p;a)x^{p-1}u^r
       =B(A-x,-u)-B(A,-u).
\end{equation}
The poles of the individual beta factors at $u=0$ cancel in this
difference.
\end{theorem}
\begin{proof}
Initially take $\Re u<0$.  The beta integral and
\eqref{harm:eq:mellinF} at $s=1$ give
\[
 F(1,u;a)=B(A,-u)=-\frac{R_a(u)}u,
 \qquad F(1,u;0)=-\frac1u.
\]
The difference of their series is normally convergent also for $u$
near zero: its $n$th term is $O(n^{|u|-2})$ on compact $a$-sets.
Consequently
\[
 \sum_{n\geq1}\prod_{j<n}(1+u/j)
       \left(\frac1n-\frac1{n+a}\right)
          =\frac{R_a(u)-1}{u}
\]
extends holomorphically through $u=0$.  Compare coefficients to
obtain \eqref{harm:eq:primitive}.  Normal convergence permits
$p-1$ derivatives in $a$, giving \eqref{harm:eq:positive}.

For an independent verification of the bivariate formula, sum
\eqref{harm:eq:mellinE} first over $p\geq2$ and then over $r$.
The result is the convergent integral
\[
 \int_0^1x_0^{A-1}(x_0^{-x}-1)(1-x_0)^{-1-u}\,dx_0.
\]
For $|x|$ sufficiently small and $|u|<1/2$, it is integrable at
both endpoints.  In the subdomain $\Re u<0$ it is a difference of
two beta integrals.  Analytic continuation across $u=0$ proves
\eqref{harm:eq:betaG} without separating divergent integrals there.
Finally, Taylor expansion of $\log\Gamma$ proves \eqref{harm:eq:R},
and logarithmic differentiation gives \eqref{harm:eq:hrec}.
\end{proof}

Formula \eqref{harm:eq:positive} is a specialization and explicit
coefficient implementation of the established parametric beta theory
\cite{Young2026}.  Its useful feature for this manuscript is that it
simultaneously evaluates infinitely many nonlinear harmonic combinations
and gives their antiderivatives.  In particular,
$\int E_r(2;a)\,da=h_{r+1}(a)+C$.

Put $Z_k=\zeta(k,A)$ and $H=H_a^{(1)}$.  Since
$\partial_a H_a^{(k)}=kZ_{k+1}$, the first outer exponent has the
especially short all-depth expression
\begin{equation}\label{harm:eq:second}
 E_r(2;a)=\sum_{k=1}^{r+1} Z_{k+1}h_{r+1-k}(a).
\end{equation}
For example,
\begin{align}
 E_1(2;a)&=HZ_2+Z_3,\label{harm:eq:examples1}\\
 E_2(2;a)&=h_2(a)Z_2+HZ_3+Z_4,\\
 E_1(3;a)&=HZ_3+\frac32Z_4-\frac12Z_2^2,\\
 E_2(3;a)&=h_2(a)Z_3+\frac32HZ_4+2Z_5
                   -\frac12HZ_2^2-Z_2Z_3.\label{harm:eq:examples4}
\end{align}
At $a=0$, \eqref{harm:eq:second} recovers
$\zeta(2,\{1\}^r)=\zeta(r+2)$; the next two examples give
$\zeta(3,1)=\zeta(4)/4$ and
$\zeta(3,1,1)=2\zeta(5)-\zeta(2)\zeta(3)$.

\subsection{An all-depth cyclotomic projection}

Use the convention
\[
 \operatorname{Li}_{p,\{1\}^r}(z_0,z_1,\ldots,z_r)
 =\sum_{m_0>m_1>\cdots>m_r\geq1}
   \frac{z_0^{m_0}z_1^{m_1}\cdots z_r^{m_r}}
        {m_0^p m_1\cdots m_r}.
\]
For $p\geq2$ this series is absolutely convergent when all the
$z_j$ are roots of unity.

\begin{corollary}[Explicit reduction of a cyclotomic average]
\label{harm:cor:cyclotomic}
Let $q\geq1$, $1\leq h\leq q$, $\omega=e^{2\pi i/q}$, and
$a=h/q-1$.  For every $p\geq2$ and $r\geq0$,
\begin{equation}\label{harm:eq:projection}
 \sum_{j_0,\ldots,j_r=0}^{q-1}\omega^{-hj_0}
 \operatorname{Li}_{p,\{1\}^r}
       (\omega^{j_0},\ldots,\omega^{j_r})
 =q^{1-p}\frac{(-1)^p}{(p-1)!}
            h_{r+1}^{(p-1)}(h/q-1).
\end{equation}
Every value on the right reduces to depth-one polylogarithms at
$q$th roots of unity and logarithms.  This is a reduction of the
specified average; it makes no assertion that its individual colored
summands all reduce to depth one.
\end{corollary}
\begin{proof}
In \eqref{harm:eq:E}, set $m_0=q(n-1)+h$ and $m_i=qk_i$ for the
indices $n-1\geq k_1>\cdots>k_r\geq1$ in $e_r(n-1)$.  This gives
\[
 E_r(p;h/q-1)=q^{p+r}
 \sum_{\substack{m_0>\cdots>m_r\geq1\\
                  m_0\equiv h\ (q),\ m_i\equiv0\ (q),\ i\geq1}}
 \frac1{m_0^p m_1\cdots m_r}.
\]
Insert the finite Fourier filters for the $r+1$ congruences and
apply \eqref{harm:eq:positive}.  To make the depth-one reduction
explicit, one has, for $k\geq2$,
\begin{align}
 \zeta(k,h/q)
 &=q^{k-1}\sum_{j=0}^{q-1}\omega^{-hj}
                       \operatorname{Li}_k(\omega^j),
                       \label{harm:eq:fourierZ}\\
 H_{h/q-1}^{(1)}
 &=-\log q+\sum_{j=1}^{q-1}\omega^{-hj}\log(1-\omega^j).
                       \label{harm:eq:fourierH}
\end{align}
The first identity follows by grouping a Dirichlet series into residue
classes.  Expand the same identity at spectral argument $1$ to obtain
the second, using $\operatorname{Li}_1(z)=-\log(1-z)$ with principal
logarithms.  Complex conjugate terms have conjugate logarithms, so
\eqref{harm:eq:fourierH} is real.  Formula \eqref{harm:eq:Hcomplex}
and the recurrence \eqref{harm:eq:hrec} finish the reduction.
\end{proof}

For a concrete nonlinear family, let $L=\log2$.  Specializing to
$a=-1/2$ gives
\begin{align}
 \sum_{n\geq0}\frac{H_n}{(2n+1)^2}
 &=\frac{7\zeta(3)-\pi^2L}{4},\label{harm:eq:odd1}\\
 \sum_{n\geq0}\frac{H_n^2-H_n^{(2)}}{(2n+1)^2}
 &=\frac{\pi^2L^2}{2}-7L\zeta(3)+\frac{\pi^4}{24},
                       \label{harm:eq:odd2}\\
 \sum_{n\geq0}\frac{H_n^3-3H_nH_n^{(2)}+2H_n^{(3)}}{(2n+1)^2}
 &=\frac32\bigl[31\zeta(5)-13\zeta(2)\zeta(3)-15L\zeta(4)
             \notag\\
 &\hspace{1.2cm}+14L^2\zeta(3)-4L^3\zeta(2)\bigr].
                       \label{harm:eq:odd3}
\end{align}
All three are instances of one identity, valid as a germ at $u=0$:
\begin{equation}\label{harm:eq:oddall}
 \sum_{r\geq0}u^r\sum_{n\geq0}\frac{e_r(n)}{(2n+1)^2}
 =\frac{\sqrt\pi\,\Gamma(1-u)}{4u\,\Gamma(1/2-u)}
       \bigl[\psi(1/2)-\psi(1/2-u)\bigr].
\end{equation}
This yields arbitrary degree in the harmonic numbers, with no
integer-relation search.  Assigning weight $k$ to $H_a^{(k)}$ and
$\zeta(k,A)$ makes \eqref{harm:eq:positive} homogeneous of weight
$p+r$; the rational Fourier reduction preserves this weight.

\subsection{Meromorphic continuation with a uniform remainder}

The same family gives exact spectral data beyond the half-plane where
its Dirichlet series converges.  Define polynomials $b_k(u,a)$ by
\begin{equation}\label{harm:eq:bdef}
 \phi(t,u,a):=e^{-At}
       \left(\frac{t}{1-e^{-t}}\right)^{1+u}
       =\sum_{k\geq0}b_k(u,a)t^k.
\end{equation}
The branch at $t=0$ is the analytic one taking value $1$.  Equivalently,
in the standard generalized Bernoulli-polynomial notation,
\[
 b_k(u,a)=\frac{B_k^{(1+u)}(u-a)}{k!}.
\]
It is enough to use the defining series; no properties of generalized
Bernoulli polynomials are assumed.  In particular,
\begin{align}
 b_0&=1,\\
 b_1&=\frac{u-2a-1}{2},\label{harm:eq:b1}\\
 b_2&=\frac{12a^2-12au+12a+3u^2-7u+2}{24},\label{harm:eq:b2}\\
 b_3&=\frac{(u-2a-1)(4a^2-4au+4a+u^2-3u)}{48}.
                       \label{harm:eq:b3}
\end{align}
Each $b_k$ is a rational polynomial in $a,u$, of degree at most $k$
in each variable.

\begin{lemma}[Uniform Mellin subtraction]\label{harm:lem:subtract}
Fix a compact subset of $\Re a>-1$ and a sufficiently small disc of
$u$-values around zero.  For every integer $K\geq1$,
\begin{equation}\label{harm:eq:subtraction}
 \Gamma(s)F(s,u;a)
   =\sum_{k=0}^{K-1}\frac{b_k(u,a)}{s-1-u+k}
       +J_K(s,u;a),
\end{equation}
where $J_K$ is jointly holomorphic on
$\Re s>1+\Re u-K$.  Therefore $F$ extends meromorphically to all
$s$, locally in $u$, and $E_r(s;a)=[u^r]F(s,u;a)$ has possible
poles only at $s=1,0,-1,-2,\ldots$.
\end{lemma}
\begin{proof}
Split \eqref{harm:eq:mellinF} at $t=1$, and define
\begin{align*}
 J_K(s,u;a)={}&\int_0^1 t^{s-2-u}
       \left(\phi(t,u,a)-\sum_{k=0}^{K-1}b_k(u,a)t^k\right)dt\\
 &+\int_1^{\infty}t^{s-1}e^{-At}
                  (1-e^{-t})^{-1-u}\,dt.
\end{align*}
The parenthesized remainder is $O(t^K)$ uniformly on the fixed
parameter sets.  Every derivative in $s$ or $u$ introduces only a
fixed power of $\log t$ near zero, so it remains integrable on compact
subsets of $\Re s>1+\Re u-K$.  The second integral and all its
parameter derivatives decay exponentially at infinity because
$\Re A$ stays positive.  Differentiation under both integrals proves
joint holomorphy.  Integrating each subtracted monomial gives
\eqref{harm:eq:subtraction}.  Increasing $K$ continues the same
function, since all formulas agree on the original convergence
domain.  Finally, $1/\Gamma(s)$ is entire; coefficient extraction
in $u$ can turn the simple moving denominators into higher poles
at the stated integers, and introduces no others.
\end{proof}

\subsection{The complete principal part and finite part at \texorpdfstring{$s=1$}{s=1}}

Write
\begin{equation}\label{harm:eq:c}
 C(u)=\frac1{\Gamma(1+u)}=\sum_{j\geq0}c_ju^j
 =\exp\left(\gamma u+
       \sum_{k\geq2}\frac{(-1)^{k+1}\zeta(k)}k u^k\right).
\end{equation}
Thus $c_0=1$, $c_1=\gamma$,
$c_2=(\gamma^2-\zeta(2))/2$, and
$c_3=(\gamma^3-3\gamma\zeta(2)+2\zeta(3))/6$.
The finite part $\operatorname{FP}_{s=s_0}$ always means the
coefficient of $(s-s_0)^0$ in the Laurent expansion of the specified
meromorphic function.

\begin{theorem}[All depths at the principal pole]\label{harm:thm:one}
For $r\geq0$ and $\Re a>-1$, as $\varepsilon\to0$,
\begin{equation}\label{harm:eq:one}
 \boxed{\quad
 E_r(1+\varepsilon;a)
 =\sum_{j=0}^{r}\frac{c_j}{\varepsilon^{r+1-j}}
       +c_{r+1}-h_{r+1}(a)+O(\varepsilon).
 \quad}
\end{equation}
In particular, all principal coefficients are independent of the
shift, while the complete shift dependence of the finite part is
the single complete harmonic polynomial $-h_{r+1}(a)$.
\end{theorem}
\begin{proof}
First take $a=0$.  With $K=1$, the singular term in
\eqref{harm:eq:subtraction} near $(s,u)=(1,0)$ is $1/(s-1-u)$.
Consequently
\[
 F(s,u;0)=\frac{1}{\Gamma(1+u)(s-1-u)}+Q(s,u),
\]
with $Q$ jointly holomorphic there: replacing $1/\Gamma(s)$ by
$1/\Gamma(1+u)$ in the singular numerator changes the expression
by a holomorphic divided difference.  For $s=1$ and small
$\Re u<0$, the beta integral gives $F(1,u;0)=-1/u$, and hence
\[
 Q(1,u)=\frac{C(u)-1}{u}.
\]
The right side extends holomorphically to $u=0$.  Extracting
$[u^r]$ from $C(u)/(\varepsilon-u)$ gives exactly
$\sum_{j=0}^{r}c_j\varepsilon^{-r-1+j}$, and extracting it from
$Q(1,u)$ gives $c_{r+1}$.

For arbitrary $a$, the difference
\[
 E_r(s;a)-E_r(s;0)
   =\sum_{n\geq1}e_r(n-1)
               \bigl((n+a)^{-s}-n^{-s}\bigr)
\]
is normally convergent for $\Re s>0$, locally uniformly in $a$.
At $s=1$ it equals $-h_{r+1}(a)$ by
\eqref{harm:eq:primitive}.  This leaves the principal part unchanged
and gives the stated finite part.
\end{proof}

For example,
\begin{align}
 \operatorname{FP}_{s=1}E_0(s;a)&=-\psi(A),\\
 \operatorname{FP}_{s=1}E_1(s;a)
 &=\frac{\gamma^2-\zeta(2)-(H_a^{(1)})^2-H_a^{(2)}}2,\\
 \operatorname{FP}_{s=1}E_2(s;a)
 &=\frac{\gamma^3-3\gamma\zeta(2)+2\zeta(3)}6-h_3(a).
\end{align}
The first line is the usual zeroth generalized Stieltjes constant.
For higher depths, the corresponding zeroth constants are just as
explicit.  This does not imply such a reduction for all higher
Laurent coefficients.

\subsection{Every nonpositive integer: a finite exact coefficient rule}

\begin{theorem}[Principal and finite Laurent data at nonpositive integers]
\label{harm:thm:negative}
Let $m,r\geq0$, and define
\begin{equation}\label{harm:eq:A}
 \mathcal A_m(u,a)=\frac{b_{m+1}(u,a)}{\Gamma(u-m)}
                =\sum_{d\geq1}A_{m,d}(a)u^d.
\end{equation}
Then
\begin{equation}\label{harm:eq:negative}
 \boxed{\quad
 E_r(-m+\varepsilon;a)
  =\sum_{h=0}^{r}A_{m,r+1-h}(a)\varepsilon^{-h}
        +O(\varepsilon).
 \quad}
\end{equation}
In particular,
\begin{equation}\label{harm:eq:FPnegative}
 \operatorname{FP}_{s=-m}E_r(s;a)
       =[u^{r+1}]\frac{b_{m+1}(u,a)}{\Gamma(u-m)}.
\end{equation}
All coefficients displayed in \eqref{harm:eq:negative} belong to
$\mathbb Q[a,\gamma,\zeta(2),\ldots,\zeta(r)]$, with the obvious
interpretation for $r=0,1$.  The highest possible pole coefficient is
\begin{equation}\label{harm:eq:leadingpole}
 A_{m,1}(a)=\zeta(-m,A).
\end{equation}
\end{theorem}
\begin{proof}
Choose $K>m+1$ in Lemma~\ref{harm:lem:subtract}.  Near
$(s,u)=(-m,0)$, precisely one of its denominators can vanish:
the one with $k=m+1$.  Thus, putting $s=-m+\varepsilon$,
\[
 F(-m+\varepsilon,u;a)
  =\frac{b_{m+1}(u,a)}{\Gamma(-m+\varepsilon)(\varepsilon-u)}
       +\frac{Q_m(\varepsilon,u;a)}{\Gamma(-m+\varepsilon)},
\]
where $Q_m$ is jointly holomorphic.  Since $1/\Gamma(-m+\varepsilon)$
has a simple zero at $\varepsilon=0$, the second term contributes
$O(\varepsilon)$ after extracting any fixed coefficient in $u$.
Write
\[
 b_{m+1}(u,a)=\sum_{j\geq0}b_{m+1,j}(a)u^j,
 \qquad \frac1{\Gamma(-m+\varepsilon)}
              =\sum_{\ell\geq1}g_{m,\ell}\varepsilon^\ell.
\]
Coefficient extraction from the singular term gives
\[
 \left(\sum_{\ell\geq1}g_{m,\ell}\varepsilon^\ell\right)
       \sum_{j=0}^r b_{m+1,j}(a)\varepsilon^{-r+j-1}.
\]
Its coefficient at $\varepsilon^{-h}$ is
$\sum_{j=0}^{r-h}b_{m+1,j}(a)g_{m,r+1-h-j}$, which is exactly
$[u^{r+1-h}]\mathcal A_m(u,a)$.  This proves
\eqref{harm:eq:negative}.

The reciprocal-gamma factor has the explicit harmonic expansion
\begin{equation}\label{harm:eq:recipgamma}
 \frac1{\Gamma(u-m)}=(-1)^m m!\,u
 \exp\left((\gamma-H_m)u+
   \sum_{k\geq2}\frac{(-1)^{k+1}\zeta(k)-H_m^{(k)}}k u^k\right).
\end{equation}
Indeed, multiply $1/\Gamma(1+u)$ by
$u\prod_{j=1}^m(u-j)$ and expand its logarithm after removing
$(-1)^m m!u$.  The finite coefficient rule now proves the claimed
polynomial ring.  Finally,
$b_{m+1}(0,a)=B_{m+1}(-a)/(m+1)!$ and
$g_{m,1}=(-1)^m m!$.  Bernoulli reflection and
$\zeta(-m,A)=-B_{m+1}(A)/(m+1)$ give
\eqref{harm:eq:leadingpole}.
\end{proof}

For a generic shift, the pole at $s=-m$ has order $r$ when $r\geq1$.
At a zero of $\zeta(-m,A)$ its order drops.  Thus the familiar
absence of negative even poles for the depth-one unshifted harmonic
zeta function is part of an all-depth cancellation rule: at $a=0$
and positive even $m$, the pole order is at most $r-1$.  The formula
determines whether further cancellations occur.

No divergent series is assigned an ordinary sum here.  The values in
\eqref{harm:eq:FPnegative} are finite Laurent coefficients of a
specified meromorphic continuation.  Nor can they be computed by
first setting $s=-m$ in $F$ at a nonzero $u$ and then differentiating:
the moving divisor $\varepsilon-u=0$ makes those operations
incompatible at their intersection.

\subsection{Small exact tables and differentiation of finite parts}

At $s=0$, \eqref{harm:eq:b1} makes the whole finite-part family
particularly simple.  With $c_{-1}=0$,
\begin{equation}\label{harm:eq:zeroall}
 \operatorname{FP}_{s=0}E_r(s;a)
       =\frac12c_{r-1}-\left(a+\frac12\right)c_r.
\end{equation}
More generally every finite part at $s=-m$ is a polynomial in $a$
of degree at most $m+1$.  The first unshifted examples are shown in
Table~\ref{harm:tab:FP}; only finite parts, not complete Laurent
series, are tabulated.

\begin{table}[htbp]
\centering
\small
\renewcommand{\arraystretch}{1.6}
\begin{tabular}{c|ccc}
 & $s=0$ & $s=-1$ & $s=-2$\\\hline
 $r=0$ & $-\tfrac12$ & $-\tfrac1{12}$ & $0$\\
 $r=1$ & $\tfrac{1-\gamma}{2}$ & $\tfrac{9-2\gamma}{24}$ & $\tfrac18$\\
 $r=2$ & $\tfrac{-\gamma^2+2\gamma+\zeta(2)}4$
       & $\tfrac{-\gamma^2+9\gamma+\zeta(2)-10}{24}$
       & $\tfrac{6\gamma-17}{48}$
\end{tabular}
\caption{Exact finite parts of $E_r(s;0)$ from the single
coefficient rule \eqref{harm:eq:FPnegative}.}
\label{harm:tab:FP}
\end{table}

For example, the complete displayed portion at a point where the
highest possible pole cancels is
\begin{align}
 E_3(-2+\varepsilon;0)
  ={}&\frac1{8\varepsilon^2}
       +\frac{6\gamma-17}{48\varepsilon}
       +\frac{3\gamma^2-17\gamma-3\zeta(2)+17}{48}
       +O(\varepsilon).\label{harm:eq:cancelledexample}
\end{align}

There is a necessary residue correction when finite parts are
differentiated.  Denote
\[
 C_{m,r}(a)=\operatorname{FP}_{s=-m}E_r(s;a)\quad(m\geq0),
 \qquad C_{-1,r}(a)=\operatorname{FP}_{s=1}E_r(s;a),
\]
and let $R_{m,r}(a)$ denote the residue at $s=-m$, also for $m=-1$.

\begin{corollary}[Differentiation with its residue term]
\label{harm:cor:residue}
For $m\geq0$,
\begin{equation}\label{harm:eq:residuecorrection}
 C_{m,r}'(a)=mC_{m-1,r}(a)-R_{m-1,r}(a).
\end{equation}
In particular $C_{0,r}'(a)=-c_r$.
\end{corollary}
\begin{proof}
Normal convergence first gives
$\partial_a E_r(s;a)=-sE_r(s+1;a)$ for $\Re s>1$;
Lemma~\ref{harm:lem:subtract} continues the identity meromorphically.
Put $s=-m+\varepsilon$.  The constant term of
$(m-\varepsilon)E_r(1-m+\varepsilon;a)$ is
$mC_{m-1,r}-R_{m-1,r}$, even if the latter function has higher
order poles.  Taking the constant coefficient commutes with
$a$-differentiation by joint meromorphy on compact shift domains.
At $m=0$, Theorem~\ref{harm:thm:one} gives $R_{-1,r}=c_r$,
in agreement with \eqref{harm:eq:zeroall}.
\end{proof}

There is also a direct coefficient-level verification.  Since
$\partial_a b_{m+1}(u,a)=-b_m(u,a)$ and
$\Gamma(u-m+1)=(u-m)\Gamma(u-m)$,
\[
 \partial_a\mathcal A_m(u,a)=(m-u)\mathcal A_{m-1}(u,a),
 \qquad m\geq1.
\]
Taking the coefficient of $u^{r+1}$ gives exactly
\eqref{harm:eq:residuecorrection}; the case $m=0$ is
\eqref{harm:eq:zeroall}.

Dropping the residue in \eqref{harm:eq:residuecorrection} would,
already at $m=0$, incorrectly predict a constant function of $a$.
This is a correction to a tempting manipulation, rather than an
assertion that the inspected manuscript makes that error.

\subsection{Higher Stieltjes coefficients and the precise scope of reduction}

Define the higher shifted height-one constants by
\begin{equation}\label{harm:eq:Stieltjes}
 E_r(1+\varepsilon;a)
 =\sum_{j=0}^{r}c_j\varepsilon^{-r-1+j}
     +\sum_{\ell\geq0}\frac{(-1)^\ell}{\ell!}
                   \gamma_{r,\ell}(a)\varepsilon^\ell.
\end{equation}
For $r=0$, these are the usual generalized Stieltjes constants
$\gamma_{0,\ell}(a)=\gamma_\ell(A)$.
Theorem~\ref{harm:thm:one} evaluates
$\gamma_{r,0}(a)=c_{r+1}-h_{r+1}(a)$.
For real $a,b>-1$ and every $\ell\geq0$, normal convergence of the
shifted difference gives the exact, absolutely convergent identity
\begin{equation}\label{harm:eq:shiftStieltjes}
 \gamma_{r,\ell}(a)-\gamma_{r,\ell}(b)
  =\sum_{n\geq1}e_r(n-1)
       \left(\frac{\log^\ell(n+a)}{n+a}
                  -\frac{\log^\ell(n+b)}{n+b}\right).
\end{equation}
The summand is $O(n^{-2}(1+\log^{r+\ell}n))$.  At $\ell=0$
this reduces to $h_{r+1}(b)-h_{r+1}(a)$.  At higher $\ell$ it
specifies the shift dependence without assuming a finite reduction
of the individual constants.

One also obtains the all-depth extension of the manuscript's
Stieltjes differentiation ladder.  Let
$e_i(k)=[v^i]\prod_{j=1}^k(1+v/j)$ and use $\partial_s$ for
spectral differentiation.  For $k\geq1$,
\begin{equation}\label{harm:eq:Stieltjesderivative}
 \frac{d^k}{da^k}\gamma_{r,\ell}(a)
 =(-1)^{\ell+k}k!\ell!
      \sum_{i+j=\ell}e_i(k)
             \frac{\partial_s^j E_r(k+1;a)}{j!}.
\end{equation}
Indeed, $\partial_a^kE_r(s;a)=(-1)^k(s)_kE_r(s+k;a)$;
the principal part at $s=1$ is independent of $a$, and coefficient
comparison at $s=1$ proves the formula.  When $\ell=0$ the right
side is fully evaluated by Theorem~\ref{harm:thm:beta}.  For
$\ell>0$ it involves genuine spectral derivatives.

For completeness these derivatives have exact logarithmic Mellin
representations.  Differentiating \eqref{harm:eq:mellinE} yields
\begin{align}
 &\frac1{r!}\int_0^1
   \frac{x^a(-\log x)^{s-1}[-\log(1-x)]^r}{1-x}
                     \log^\ell(-\log x)\,dx\notag\\
 &\hspace{1cm}
 =\sum_{j=0}^{\ell}\binom{\ell}{j}
      \Gamma^{(\ell-j)}(s)\,\partial_s^jE_r(s;a),
      \qquad \Re s>1.\label{harm:eq:logspectral}
\end{align}
The endpoint estimates in Proposition~\ref{harm:prop:mellin} remain
integrable after any fixed number of logarithmic factors, justifying
all differentiations.  Gamma derivatives here can equivalently be
written using complete Bell polynomials in polygamma functions.

The exact polynomial conclusion of
Theorem~\ref{harm:thm:negative} stops at the constant Laurent
coefficient.  The term $Q_m/\Gamma(s)$ in its proof starts at
order $s+m$ and can contribute new constants from that point
onward.  At depth zero these already include $\zeta'(-m,A)$,
and at $m=0$ Lerch's identity gives the log-gamma function.
There is therefore no basis for discarding the analytic remainder
when computing positive spectral derivatives.

The next concrete questions are to identify which linear combinations
of these higher coefficients eliminate the analytic remainder, to
extend the cyclotomic projection to harmonic patterns other than
$\{1\}^r$, and to compare the outer-shift constants
\eqref{harm:eq:Stieltjes} systematically with the literature's
simultaneously shifted harmonic Hurwitz constants.  Those extensions
are left open here.
